\paragraph{Section 9 fragment statement.}
\paragraph{Motivating conjecture from the repair brief, verbatim.}
(NOT YET PROVED.)  Under Assumptions~\ref{ass:sampling}--\ref{ass:regular}, the weak-duality inequalities in Theorem~\ref{thm:k-sharp-dual-confirmed} are equalities:
\[
L_x(P)=D_x(h_x),\qquad U_x(P)=-D_x(-h_x)
\]
for almost every \(x\), and hence the integrated dual values equal \(L_\rho(P)\) and \(U_\rho(P)\).  The conjecture does not assert endpoint attainment; it asserts strong duality for the bounded Borel strict-indicator cost \(h_x\) under the uniform-margin and Spearman-band constraints.

\paragraph{Original proposal claim recorded in the setup.}
(NOT YET PROVED.)  Under Assumptions~\ref{ass:sampling}--\ref{ass:band}, the sharp identified set for \(\theta=\Pr(Y_1>Y_0)\) is exactly the closed interval \([L_\rho(P),U_\rho(P)]\).  For each \(x\), the conditional value set \(I_x=\{\int h_x\,d\gamma:\gamma\in\Gamma_x(\rho_-,\rho_+)\}\) is the closed interval \([L_x(P),U_x(P)]\), the lower endpoint is attained by a measurable copula selector \(x\mapsto\gamma_x^L\) and equals
\[
\sup_{\alpha,\beta,\lambda_-,\lambda_+}
\left\{\int_0^1\alpha(u)\,du+\int_0^1\beta(v)\,dv+\lambda_-a_- - \lambda_+a_+\right\},
\]
where \(\lambda_-,\lambda_+\ge0\) and
\[
\alpha(u)+\beta(v)\le h_x(u,v)+(\lambda_+-\lambda_-)uv\quad\text{for all }(u,v)\in[0,1]^2,
\]
and the upper endpoint is attained by a measurable selector \(x\mapsto\gamma_x^U\) and is obtained by applying the same formula to \(-h_x\) and changing sign.

\paragraph{Headline statement.}
Under Assumptions~\ref{ass:sampling}--\ref{ass:regular}, the available assumptions imply the computable weak Spearman-band dual certificate, but they do not prove the conjectured no-gap equalities.  For \(P_X\)-almost every \(x\) and every \(g\in\{h_x,-h_x\}\), \[D_x(g)\le \inf_{\gamma\in\Gamma_x(\rho_-,\rho_+)}\int g\,d\gamma,\] where \[D_x(g)=\sup_{\lambda_-,\lambda_+\ge0,\ \alpha,\beta\in L^1([0,1])}\left\{\int_0^1\alpha(u)\,du+\int_0^1\beta(v)\,dv+\lambda_-a_- -\lambda_+a_+:\alpha(u)+\beta(v)\le g(u,v)+(\lambda_+-\lambda_-)uv\ \forall(u,v)\in[0,1]^2\right\}.\] Consequently every feasible latent law obeys \[\int D_x(h_x)\,dP_X(x)\le \theta\le \int\{-D_x(-h_x)\}\,dP_X(x).\] This is a partial repair only: the finite two-strata refutation search below finds no counterexample, but the reverse inequalities \(L_x(P)\le D_x(h_x)\) and \(-D_x(-h_x)\le U_x(P)\), the endpoint-attaining measurable selectors, closedness of \(I_x\), and the interpolation step needed for \(\Theta_I(P)=[L_\rho(P),U_\rho(P)]\) remain unproved under the written assumptions.

\begin{proposition}[Partial repair of Conjecture 1: weak Spearman-band dual certificate]\label{thm:k-sharp-dual-confirmed}
Under Assumptions~\ref{ass:sampling}--\ref{ass:regular}, for \(P_X\)-almost every \(x\) and every \(g\in\{h_x,-h_x\}\),
\[
D_x(g)\le
\inf_{\gamma\in\Gamma_x(\rho_-,\rho_+)}\int g\,d\gamma,
\]
where
\[
D_x(g)=\sup_{\lambda_-,\lambda_+\ge0,\ \alpha,\beta\in L^1([0,1])}
\left\{\int_0^1\alpha(u)\,du+\int_0^1\beta(v)\,dv+\lambda_-a_- -\lambda_+a_+:
\alpha(u)+\beta(v)\le g(u,v)+(\lambda_+-\lambda_-)uv\ \forall(u,v)\in[0,1]^2
\right\}.
\]
Consequently every feasible latent law obeys
\[
\int D_x(h_x)\,dP_X(x)\le \theta\le
\int\{-D_x(-h_x)\}\,dP_X(x).
\]
The proposition is only a weak-duality certificate.  It does not establish the reverse inequalities \(L_x(P)\le D_x(h_x)\) or \(-D_x(-h_x)\le U_x(P)\), endpoint attainment, measurable argmax selectors, or the closed sharp-identified-set equality \(\Theta_I(P)=[L_\rho(P),U_\rho(P)]\).  Thus Conjecture 1 is partial under the current assumption set.
\end{proposition}

\paragraph{Section 13 refutation attempt before proof.}
I first applied the required two-strata response-type LP and vertex-enumeration check to the smallest nontrivial rank-copula approximation.  Take \(X\) degenerate, use the two rank grid points \(u_1=1/3\) and \(u_2=2/3\), impose row and column masses \(1/2\), and write
\[
\pi^+=\begin{pmatrix}1/2&0\\0&1/2\end{pmatrix},
\qquad
\pi^-=\begin{pmatrix}0&1/2\\1/2&0\end{pmatrix}.
\]
Every feasible \(2\times2\) coupling is
\[
\pi_t=t\pi^+ +(1-t)\pi^-,\qquad 0\le t\le1,
\]
and its Spearman moment is
\[
m(t)=\sum_{\ell,m}u_\ell u_m\pi_{t,\ell m}
=\frac{2}{9}+\frac{t}{18}.
\]
Thus the moment band restricts \(t\) to the explicit interval
\[
T(a_-,a_+)=
\left[\max\{0,18a_- -4\},\min\{1,18a_+ -4\}\right].
\]
The threshold-zero benefit kernels compatible with increasing quantiles are the staircase binary matrices.  Apart from the all-zero and all-one cases, the four nonconstant patterns are
\[
\begin{pmatrix}0&0\\1&0\end{pmatrix},\quad
\begin{pmatrix}0&0\\1&1\end{pmatrix},\quad
\begin{pmatrix}1&0\\1&0\end{pmatrix},\quad
\begin{pmatrix}1&0\\1&1\end{pmatrix}.
\]
Their primal values under \(\pi_t\) are, respectively,
\[
\frac{1-t}{2},\qquad
\frac{1}{2},\qquad
\frac{1}{2},\qquad
\frac{1+t}{2}.
\]
Optimizing these affine objectives over \(T(a_-,a_+)\) gives an endpoint at a face of the finite LP.  The finite-dimensional dual has variables for the two row constraints, the two column constraints, and the two Spearman inequalities; ordinary finite LP strong duality gives the same face value in every staircase case.  Therefore no vertex or dual certificate in the two-strata search violates the proposed equality \(L_x(P)=D_x(h_x)\) or \(U_x(P)=-D_x(-h_x)\).  This refutation attempt also checks the Frechet equality faces in the smallest grid: saturation of \(t=\min T(a_-,a_+)\) or \(t=\max T(a_-,a_+)\) selects the finite dual face rather than producing a gap.  The finite search therefore fails to refute the bounded-Borel strong-duality statement.  It also cannot prove it, because ordinary finite LP duality does not control the continuous bounded-Borel marginal problem with an added Spearman moment interval.  Nor can it confirm the original sharp identified-set claim, because every finite LP has endpoint-attaining optimizers and so does not test the continuous measurable-selector and residual-interpolation problem flagged by the reviewer.

\paragraph{Section 10 interpretation.}
The repair separates what is proved from what would require a separate moment-constrained bounded-Borel transport theorem.  The functions \(\alpha\) and \(\beta\) are the Kellerer-style marginal potentials for the two uniform rank margins, while \(\lambda_-\) and \(\lambda_+\) price the lower and upper Spearman moment inequalities.  Any such feasible potential system gives a computable lower certificate for \(L_x(P)\) and, after applying the same construction to \(-h_x\), a computable upper certificate for \(U_x(P)\).  The missing step is the reverse inequality: the current assumptions do not state, and the proof below does not derive, a no-gap theorem for a bounded Borel strict-indicator cost under both marginal constraints and a Spearman moment interval.

\paragraph{Section 11 proof.}
\begin{proof}
Fix a \(P_X\)-relevant covariate value \(x\).  Assumptions~\ref{ass:sampling}--\ref{ass:continuous} identify the conditional quantile maps \(Q_1(\cdot\mid x)\) and \(Q_0(\cdot\mid x)\) from the observed law \(P\).  Hence
\[
h_x(u,v)=\one\{Q_1(u\mid x)>Q_0(v\mid x)\}
\]
is a bounded Borel function on the compact Polish space \([0,1]^2\).  Assumption~\ref{ass:band} fixes the feasible rank-copula class
\[
\Gamma_x(\rho_-,\rho_+)=
\left\{\gamma\in\Pi(\lambda,\lambda):
a_-\le \int uv\,d\gamma(u,v)\le a_+\right\},
\qquad
a_\pm=\frac{\rho_\pm+3}{12}.
\]
The class is nonempty because \(-1<\rho_-\le\rho_+<1\), the countermonotone copula has \(\int uv\,d\gamma=1/6\), the comonotone copula has \(\int uv\,d\gamma=1/3\), and mixtures attain every moment in the interval \([1/6,1/3]\).  The uniform-margin restrictions and the two moment inequalities are the only restrictions used in the duality step.

First prove weak duality.  Let \(g\in\{h_x,-h_x\}\), let \(\gamma\in\Gamma_x(\rho_-,\rho_+)\), and take any dual tuple \((\alpha,\beta,\lambda_-,\lambda_+)\) with \(\lambda_-,\lambda_+\ge0\) and
\[
\alpha(u)+\beta(v)\le g(u,v)+(\lambda_+-\lambda_-)uv
\quad\text{for all }(u,v)\in[0,1]^2.
\]
Integrating the constraint against \(\gamma\), using the two uniform margins, and adding the multiplier constants gives
\[
\begin{aligned}
&\int_0^1\alpha(u)\,du+\int_0^1\beta(v)\,dv+\lambda_-a_- -\lambda_+a_+\\
&\qquad\le
\int g\,d\gamma
+\lambda_+\left(\int uv\,d\gamma-a_+\right)
+\lambda_-\left(a_--\int uv\,d\gamma\right)
\le \int g\,d\gamma.
\end{aligned}
\]
Taking the supremum over all admissible dual tuples and then the infimum over feasible \(\gamma\) yields
\[
D_x(g)\le \inf_{\gamma\in\Gamma_x(\rho_-,\rho_+)}\int g\,d\gamma.
\]

The displayed \(D_x(g)\) is the Kellerer marginal-dual template \cite{Kellerer1984Duality} augmented by the two nonnegative Lagrange multipliers for the Spearman interval.  The preceding calculation proves the direction that follows from the written assumptions alone.  To obtain equality one would need to interchange the infimum over \(\gamma\) and the supremum over \((\lambda_-,\lambda_+,\alpha,\beta)\), or equivalently prove a moment-constrained bounded-Borel transport duality theorem for the cost \(g\) and the continuous moment \(uv\).  Assumption~\ref{ass:regular} expressly does not state such a theorem, and the present repair does not assume it.

Taking \(g=h_x\) gives the lower certificate:
\[
D_x(h_x)\le
\inf_{\gamma\in\Gamma_x(\rho_-,\rho_+)}\int h_x\,d\gamma
=L_x(P).
\]
Taking \(g=-h_x\) and multiplying by \(-1\) gives the upper certificate:
\[
U_x(P)
=\sup_{\gamma\in\Gamma_x(\rho_-,\rho_+)}\int h_x\,d\gamma
=-\inf_{\gamma\in\Gamma_x(\rho_-,\rho_+)}\int (-h_x)\,d\gamma
\le -D_x(-h_x).
\]
Now let \(x\mapsto\gamma_x\) be any feasible conditional copula selection for a latent law satisfying Assumption~\ref{ass:band}.  Applying the two pointwise weak-duality inequalities to \(\gamma_x\) yields, for \(P_X\)-almost every \(x\),
\[
D_x(h_x)\le \int h_x\,d\gamma_x\le -D_x(-h_x).
\]
Integrating over \(P_X\) gives the advertised certificate for the causal target
\[
\int D_x(h_x)\,dP_X(x)\le
\theta(\gamma_\cdot)
\le \int\{-D_x(-h_x)\}\,dP_X(x).
\]
No step in this argument proves the reverse inequalities, asserts that either infimum is attained, shows that \(I_x\) is closed, or constructs endpoint selectors \(x\mapsto\gamma_x^L,\gamma_x^U\).  To prove the original proposal's sharp identified-set equality one would still need the missing no-gap theorem together with feasible measurable selectors, or epsilon-selectors with residuals
\[
r_L(\epsilon)=\int\left\{\int h_x\,d\gamma_x^{\epsilon,L}-L_x(P)\right\}\,dP_X(x)\to0,\qquad
r_U(\epsilon)=\int\left\{U_x(P)-\int h_x\,d\gamma_x^{\epsilon,U}\right\}\,dP_X(x)\to0,
\]
plus an interpolation argument showing that all values between the integrated endpoints are feasible.  Those selector and interpolation claims are not consequences of the weak certificate proved here.  The proof therefore establishes only a computable dual certificate and leaves the strong-duality and sharpness claims unresolved.
\end{proof}

\paragraph{Section 12 sanity check.}
Let the Spearman band expand to the vacuous copula range, so \(a_-\downarrow1/6\) and \(a_+\uparrow1/3\).  Every copula with uniform margins satisfies
\[
\frac16\le \int uv\,d\gamma\le \frac13,
\]
with equality at the countermonotone and comonotone copulas.  The Spearman constraints then impose no restriction, and the multiplier-free certificate has \(\lambda_-=\lambda_+=0\).  In this corner the missing reverse direction is supplied by Kellerer's ordinary marginal duality
\[
\inf_{\gamma\in\Pi(\lambda,\lambda)}\int h_x\,d\gamma
=
\sup_{\alpha+\beta\le h_x}
\left\{\int_0^1\alpha(u)\,du+\int_0^1\beta(v)\,dv\right\},
\]
with the analogous formula for \(-h_x\).  This is the continuous-outcome Makarov/Fan--Park copula-free endpoint benchmark for the threshold-zero benefit event.  The check confirms that the weak certificate has the right form when rank information is removed, but it does not supply the moment-constrained reverse inequality, measurable selectors, or interpolation step missing from the original sharp identified-set claim.

\paragraph{Section 13 counterexample or minimality check.}
The Spearman multipliers are necessary; dropping them would compute the copula-free marginal problem rather than the band-restricted endpoint.  The same two-rank-grid calculation from the refutation attempt gives a concrete certificate.  For the staircase kernel
\[
h=\begin{pmatrix}1&0\\1&1\end{pmatrix},
\]
the objective value is \((1+t)/2\) on \(\pi_t=t\pi^+ +(1-t)\pi^-\).  Without the Spearman restriction, \(t=1\) is feasible and the grid upper value is \(1\).  Under a strict interior band with \(\max T(a_-,a_+)<1\), the correct upper value is instead
\[
\frac{1+\max T(a_-,a_+)}{2}<1.
\]
Thus a dual certificate that omits \(\lambda_-\) and \(\lambda_+\) fails to evaluate the locked band functional even at the finite-grid level.  The rank-moment restriction is not a cosmetic regularizer; it is the finite dependence information whose support-function term produces the displayed constrained dual.

The selector residuals flagged in the review are also load-bearing for the original sharp identified-set claim.  Pointwise equalities \(L_x(P)=D_x(h_x)\) and \(U_x(P)=-D_x(-h_x)\) imply only
\[
L_\rho(P)\le \theta(\gamma_\cdot)\le U_\rho(P)
\]
for every feasible measurable selection \(x\mapsto\gamma_x\), and the repaired proposition does not even establish those pointwise equalities.  They would not by themselves produce selections whose integrated values approach the endpoints, nor would they show that intermediate values between two feasible selections are themselves feasible under the conditional Spearman band.  Removing the no-gap, selector, and residual steps is exactly what turns the original sharpness theorem into the weaker weak-dual certificate proved above.
